\subsection{应用：II. 连续复函数的向量空间}

设 X 是一个非空紧空间，H 是复 Banach 空间 \(\mathcal{C}(\mathrm{X};\mathbf{C})\) 的一个线性子空间，它包含常数且分离 X 的点。诸 \(\mathcal{R}(f)\)（函数 \(f \in H\) 的实部）所成之集是一个线性子空间 \(H_{r}\)（实向量空间 \(\mathcal{C}(\mathrm{X};\mathbf{R})\) 的）；对每个 \(f \in H\)，集 \(H_{r}\) 还包含 \(\mathcal{I}(f) = \mathcal{R}(-if)\)；由此推出 \(H_{r}\) 分离 X 的点，因为关系 \(h(x) = h(y)\)（对一切 \(h \in H_{r}\)）蕴涵 \(\mathcal{R}(f(x)) = \mathcal{R}(f(y))\) 与 \(\mathcal{I}(f(x)) = \mathcal{I}(f(y))\)，从而 \(f(x) = f(y)\)（对一切 \(f \in H\)）。X 中的 \(H_{r}\)-极值点仍称为 H-极值点，其全体记作 \(\operatorname{Ch}_{\mathcal{H}}(\mathrm{X})\)，后者的闭包记作 \(\check{\operatorname{S}}_{\mathcal{H}}(\mathrm{X})\)。命题 5 与命题 7 的类似物如下：

命题 9. --- 对每个函数 \(f \in \mathcal{H}\)，\(\mathrm{Ch}_{\mathcal{H}}(\mathrm{X})\) 与使 \(|f|\) 达到其上确界的点所成之集相交。

我们可以限于 \(f\) 不是常数 0 的情形。设 \(a\) 是 \(X\) 中使 \(|f|\) 达到其上确界的一个点，并令 \(g = f / f(a)\)；那么 \(g(a) = 1\) 且 \(|g(x)| \leqslant 1\)（对一切 \(x \in X\)），于是

\[
\mathscr {R} (g (a)) = 1 \quad \text { and } \quad \mathscr {R} (g (x)) \leqslant 1 \quad \text { for   all } x \in X.
\]

把 No.~3 的命题 5 应用于 \(\mathcal{H}_r\)，存在 \(b \in \mathrm{Ch}_{\mathcal{H}}(\mathrm{X})\) 使 \(\mathcal{R}(g(x))\) 在该点达到其上确界 1，于是 \(|g(b)| = 1\)（因为 \(|g(b)| \leqslant 1\)）；由此得出 \(|f(b)| = |f(a)| \geqslant |f(x)|\)（对一切 \(x \in \mathrm{X}\)）。

命题 10. --- 设 F 是 X 的一个闭子集。下列性质等价：

\begin{enumerate}
\def\labelenumi{\alph{enumi})}
\item
  F 包含 \(\check{\mathrm{S}}_{\mathcal{H}}(\mathrm{X})\)
\item
  对每个函数 \(f \in \mathcal{H}\)，\(F\) 都与 \(X\) 中使 \(|f|\) 达到其上确界的点所成之集相交。
\item
  对每个点 \(x \in X\)，存在 X 上一个总质量为 1 的正测度 \(\mu\)，使得 \(\operatorname{Supp}(\mu) \subset \operatorname{F}\) 且 \(f(x) = \int f \, d\mu\)（对每个函数 \(f \in H\)）。
\end{enumerate}

我们来证明条件 \(a)\) 与 \(c)\) 的等价性：设 \(f = f_{1} + if_{2}\)，其中 \(f_{1}, f_{2}\) 属于 \(\mathcal{H}_{r}\)；关系 \(f(x) = \int f d\mu\) 等价于两个关系 \(f_{1}(x) = \int f_{1} d\mu\) 与 \(f_{2}(x) = \int f_{2} d\mu\)；于是对 \(\mathcal{H}_{r}\) 只须应用 No.~3 命题 7 中条件 \(a)\) 与 \(c)\) 的等价性。\(a)\) 蕴涵 \(b)\) 这一事实由命题 9 得出。我们来证明 \(b)\) 蕴涵 \(a)\)；这就是要看到：若 \(b)\) 成立，则对每个 \(h \in \mathcal{H}_{r}\)，\(F\) 都与 \(h\) 在 \(X\) 中达到其下确界的点所成之集相交。条件 \(b)\) 蕴涵 \(F\) 非空；由于 \(F\) 是紧的，存在 \(a \in F\) 使得 \(h(a) \leqslant h(y)\)（对一切 \(y \in F\)）。设 \(f \in \mathcal{H}\) 满足 \(h = \mathcal{R}(f)\)；对每个 \(\varepsilon > 0\)，函数 \(g = f - h(a) + \varepsilon\) 属于 \(\mathcal{H}\)，且

\[
\mathcal {R} (g (y)) = h (y) - h (a) + \varepsilon \geqslant \varepsilon
\]

对一切 \(y \in F\) 成立。用 c 表示 \(|g|\) 在 X 中的上确界，并令 \(b = c^{2}/2\varepsilon\)；对每个 \(y \in F\)，

\[
\left| g (y) - b \right| ^ {2} = \left| g (y) \right| ^ {2} - 2 b \mathcal {R} (g (y)) + b ^ {2} \leqslant c ^ {2} - 2 b \varepsilon + b ^ {2} = b ^ {2},
\]

换言之，函数 \(|g - b|\) 在 F 中的上确界 \(\leqslant b\)。由于 \(g - b \in \mathcal{H}\)，关于 F 的假设蕴涵 \(|g - b| \leqslant b\)，于是

\[
b ^ {2} \geqslant | g - b | ^ {2} = | g | ^ {2} - 2 b \mathcal {R} (g) + b ^ {2}
\]

从而 \(\mathcal{R}(g) \geqslant |g|^{2}/2b \geqslant 0\)；由于 \(\mathcal{R}(g) = h - h(a) + \varepsilon\) 且 \(\varepsilon > 0\) 任意，我们有 \(h \geqslant h(a)\)，而 \(h(a)\) 是 h 在 X 中的下确界，证明至此完成。

注. --- 若 f 是一个连续实函数，则使 \(|f|\) 达到其上确界的点也是函数 f 与 -f 之一达到其上确界的点。对满足 No.~3 中假设的连续实函数的向量空间 H 而言，命题 9 与命题 10 因而是命题 5 与命题 7 的平凡推论。

